% 02_background.tex
% Acronyms defined on first use, per prompts/11_paper.md style rule #5.

\section{Background}\label{sec:background}

This section introduces the Model Context Protocol (\mcp{}) and
recapitulates the isolation vocabulary we use throughout the paper.

\subsection{The Model Context Protocol}

\mcp{} (\mcp{}) is a JSON-RPC-2.0-based protocol that mediates
communication between an LLM agent (the \emph{host}) and one or
more external \emph{servers} that expose tools, resources, and
prompt templates~\cite{mcp_spec}. The protocol defines five
primitive message types:
\texttt{initialize}, \texttt{tools/list},
\texttt{tools/call}, \texttt{resources/read}, and
\texttt{prompts/get}, plus a notifications channel for
server-pushed events. Transports include stdio (for
co-located processes), HTTP+SSE, and streamable HTTP. An MCP
\emph{host} (e.g. an agent runtime) may multiplex many
connections to different servers; an MCP \emph{server} may serve
many \emph{tenants} (users, agents, or organisations).

\paragraph{Acronyms.}
\textbf{API} Application Programming Interface; \textbf{CLI}
Command-Line Interface; \textbf{CVSS} Common Vulnerability
Scoring System; \textbf{CWE} Common Weakness Enumeration;
\textbf{DoS} Denial of Service; \textbf{JSON-RPC} JSON Remote
Procedure Call; \textbf{JWT} JSON Web Token; \textbf{LLM} Large
Language Model; \textbf{LSP} Language Server Protocol;
\textbf{mTLS} Mutual Transport Layer Security; \textbf{MCP} Model
Context Protocol; \textbf{NDSS} Network and Distributed System
Security; \textbf{PII} Personally Identifiable Information;
\textbf{RPC} Remote Procedure Call; \textbf{SSE} Server-Sent
Events; \textbf{STRIDE} Spoofing, Tampering, Repudiation,
Information Disclosure, Denial of Service, Elevation of
Privilege~\cite{stride_orig}; \textbf{SAN} Subject Alternative
Name; \textbf{SSRF} Server-Side Request Forgery;
\textbf{TOCTOU} Time-of-Check Time-of-Use.

\subsection{Isolation in multi-tenant agent infrastructure}

We use \emph{isolation} in the security-engineering sense: a
tenant's data, computation, and identity are confined to that
tenant, with no observable influence on---or from---other
tenants~\cite{nist_sp_800_207}. The MCP specification does not
mandate any isolation mechanism; it simply defines the wire
format. Reference implementations (Python, TypeScript, Rust)
provide single-tenant defaults and leave multi-tenancy as an
implementation concern.

\paragraph{Why MCP's eight boundaries matter.}
We catalogue eight implicit boundaries that any multi-tenant
MCP server must enforce. Each boundary corresponds to a
trust-decision point in the protocol:

\begin{itemize}
  \item \textbf{Transport}---the byte-level channel between host
        and server (stdio, HTTP+SSE, streamable HTTP).
  \item \textbf{Session}---per-connection state on the server
        (resolved principal, caches, cancellation tokens, SSE
        queues).
  \item \textbf{Namespace}---tool, resource, and prompt naming
        and discovery.
  \item \textbf{Tool}---tool invocation and result delivery.
  \item \textbf{Resource}---file/blob/URI access via MCP
        resources.
  \item \textbf{Memory}---server-side conversation memory and
        scratchpads.
  \item \textbf{Cache}---cached tool outputs and embeddings.
  \item \textbf{Auth}---token and scope verification.
\end{itemize}

A boundary is \emph{enforced} when every observable
interactions across the boundary carries an unambiguous
tenant identity that the server can verify before acting on the
request. A boundary is \emph{broken} when an actor without a
legitimate tenant identity can either read another tenant's
state or cause another tenant to act on their behalf. We
operationalise these definitions in \S\ref{sec:threat_model}.

\subsection{Threat-model vocabulary}

We adopt STRIDE as our per-boundary threat taxonomy because it
is the de-facto vocabulary in protocol-level threat
modelling~\cite{stride_orig}. Each STRIDE letter maps cleanly
onto an MCP failure mode:

\begin{itemize}
  \item \textbf{Spoofing (S)}: impersonate a tenant or forge a
        session/transport identity.
  \item \textbf{Tampering (T)}: modify another tenant's state
        via a shared write path.
  \item \textbf{Repudiation (R)}: hide or falsify the actor of
        an attack via missing audit.
  \item \textbf{Information Disclosure (I)}: read another
        tenant's data.
  \item \textbf{Denial of Service (D)}: starve or degrade
        another tenant's service.
  \item \textbf{Elevation of Privilege (E)}: bypass authorisation
        via shadow or renamed capability.
\end{itemize}

The cross-product of eight boundaries and six STRIDE letters
yields 48 rows, all of which we enumerate in
\S\ref{sec:threat_model}.

\subsection{Defense vocabulary}

We evaluate four protocol-layer defenses (rather than OS-layer
sandboxing) because the goal is to fix MCP, not the underlying
operating system~\cite{dipta_2024_sandbox}. The four defenses
are: \emph{per-tenant tool registries} (namespace), \emph{tenant-
prefixed cache keys} (cache), \emph{resource-path
canonicalisation} (resource), and \emph{audience-bound JWTs}
(auth). Each defense is a self-contained middleware that can be
enabled or disabled per experiment, and we measure them
individually and in composition in \S\ref{sec:evaluation}.

\subsection{Why the eight boundaries matter together}

A common objection to boundary-by-boundary analysis is that
real-world attacks cross boundaries: a single prompt-injection
attack can traverse the resource boundary (the injected content
is read via \texttt{resources/read}), the auth boundary (the
injected instruction seeks elevation), and the tool boundary
(the model is tricked into calling a tool). Our defense
blueprint (\S\ref{sec:defenses}) addresses this objection by
arguing that the composition of one defense per boundary is
sufficient: each attack's exploitation path crosses
\emph{at least one} boundary whose defense is engaged. Cache
attacks cross the namespace boundary (tool routing) and the
cache boundary; the registry defence catches the former and
the cache defence catches the latter. Prompt-injection attacks
cross the resource boundary and the auth boundary; the
canonicalisation and JWT defences catch them in series.

A second common objection is that boundary catalogues are
\emph{not falsifiable} --- that any failure can be re-described
in terms of any boundary. We address this by insisting that
each row in our STRIDE table has a \emph{mechanism}: a
concrete code path or protocol message whose presence or
absence determines whether the boundary is enforced. The
per-tenant tool registry, for example, is a single middleware
line; if it is absent, the row \texttt{A-NSP-001} (tool-name
squatting) is exploitable; if it is present, the row is not.
The mechanism makes the catalogue falsifiable and the defenses
testable.

\subsection{Capability-based security as the underlying model}

The defense blueprint in \S\ref{sec:blueprint} is consistent with
\emph{capability-based security}: every action requires an
unforgeable capability that names both the principal and the
target. Capsicum (FreeBSD 9.0, 2012) is the historical
precedent; the audience-bound JWT in our auth defense is the
modern analogue. The key insight is that \emph{ambient authority}
--- permissions inherited from the process identity rather than
explicitly granted per action --- is the root cause of most
multi-tenant security failures. MCP's defaults inherit ambient
authority from the stdio worker (transport boundary), the
in-memory cache (cache boundary), and the shared tool registry
(namespace boundary); our defenses replace each ambient
authority with an explicit capability.

\subsection{Why protocol-layer defenses, not OS-layer sandboxes}

Dipta et al.~(2024) demonstrated that side channels survive
logical sandboxing~\cite{dipta_2024_sandbox}: their dynamic
frequency-based fingerprinting attack defeats gVisor,
Firecracker, SGX, and SEV. This is a \emph{positive} result for
their attack and a \emph{negative} result for OS-layer
sandboxing as a sole defense. Our position is that MCP-layer
defenses and OS-layer sandboxes are \emph{complementary}: the
MCP-layer defenses close the protocol-level leakage paths (the
RQ-4 result), and the OS-layer sandbox contains the residual
silicon-level leakage (Dipta et al.'s threat model). A
production MCP deployment should run the secure reference
server inside a Firecracker microVM and engage all four
protocol-layer defenses; the combination is what the MCP
specification should recommend.